\chapter{Netzwerkkonfiguration in einer VM}
\label{kap:netzwerk}

\begin{aufgabe}[NAT-Konfiguration]
Stelle die Netzwerkverbindung der VM auf \enquote{NAT} und verbinde dich mit dem Netzwerk. Füge einen Screenshot ein, auf dem die Einstellungen ersichtlich sind und einen Screenshot auf dem ersichtlich ist, welche IP-Adresse dein Host-PC und dein Guest-PC besitzen. Führe einen Ping vom Guest-PC zu deinem Host-PC aus und umgekehrt. Hat der Ping funktioniert? Begründe. Versuche einen Zugriff auf eine Seite im Internet. Hat der Zugriff funktioniert? Begründe.
\end{aufgabe}

% TODO Antwort

\begin{aufgabe}[Die drei Netzwerkmodi]
Erkläre die drei Netzwerkeinstellungen \enquote{NAT}, \enquote{Bridged} und \enquote{Host only} aus dem Handbuch Kap. 6.2ff. Beschreibe jeden Begriff kurz, und gib an, welche Verbindungen zwischen Guest, Host und Internet funktionieren werden. Erstelle weiters zu jeder Konfiguration eine Skizze mit Internet, Host und Guest auf der die Netzwerktopologie ersichtlich ist.
\end{aufgabe}

% TODO NAT, Bridged, Host-only je mit Tabelle und Skizze

\begin{aufgabe}[Bridged-Konfiguration]
Stelle die Netzwerkverbindung der VM auf \enquote{Bridged} und verbinde dich mit dem Netzwerk. Gib der virtuellen Maschine die IP-Adresse des Rechners + 50, Standardgateway ist 172.16.110.254, DNS ist 172.16.1.231. Erstelle einen Screenshot der IP-Konfiguration von Host-PC UND Guest-PC. Führe einen Ping vom Guest-PC zu deinem Host-PC aus und umgekehrt. Hat der Ping funktioniert? Begründe. Versuche aus dem Guest einen Zugriff auf eine Seite im Internet. Hat der Zugriff funktioniert? Begründe.
\end{aufgabe}

% TODO Antwort
